\section{Purpose}

\textbf{[STIPULATION]} This document closes the gap between half-integrality
on the torus (A263) and chiral projection into flat space (A090, A264). The
core: a torus has an inside and an outside with topologically different coupling
directions. After decompactification these appear as two chiral states -- and
the projection is not symmetric.

\section{Half-Integrality and Torus Geometry}

\textbf{[B]} Spin-$\tfrac12$ follows from the winding $(n_\theta,n_\phi)=(1,1)$
on $T^4$ (A263): two rotations $(720°)$ bring the spinor back to its initial
state. This is pure topology.

\textbf{[B]} A torus with winding $(1,1)$ has an intrinsic distinction: inside
and outside. For an embedded torus in $\mathbb{R}^3$: the Gaussian curvature is
positive on the outside ($K>0$) and negative on the inside ($K<0$). For the
flat identification torus $T^4$, $K=0$ everywhere -- but the poloidal/toroidal
circulation structure nevertheless defines a topological inside/outside
distinction through the orientation of the winding.

\textbf{[B]} The field direction is topologically fixed by the winding geometry:
the electric field vector points inward on the outside (for negative charge) and
outward (for positive charge). This direction is not freely choosable -- it
follows from the quantised winding number $\Phi=n\cdot h/e$.

\section{Chiral Projection}

\textbf{[B]} The decompactification $T^4\to\mathbb{R}^3\times\mathbb{R}_t$
(type-I projection, A090) is not symmetric with respect to inside and outside.
The compact torus carries both sides simultaneously; the decompactified flat
space sees only one projection.

\textbf{[B]} What remains after projection:
\begin{itemize}[leftmargin=2em,itemsep=2pt,topsep=2pt]
  \item The \emph{outer state} projects to left-handed coupling
    (positive rotation phase $+\xi/2$).
  \item The \emph{inner state} projects to right-handed coupling
    (negative rotation phase $-\xi/2$).
\end{itemize}
The projection selects one side -- this is the geometric origin of chiral
asymmetry.

\textbf{[B]} \textbf{Why this explains the chiral structure.} In the compact
$T^4$, both sides couple with equal strength -- there is no preferred direction.
After decompactification, only the outer perspective exists as an observable
state. The inner state is not directly accessible through the type-I projection
-- it is the ``other branch'' that is lost in the dimensional reduction (A090,
type-II loss).

\section{Connection to the Weak Interaction}

\textbf{[B]} \textbf{Fourier projector and $g_R=0$.} The Hilbert space on $T^4$
decomposes into sectors by the winding index $m$. The physical projector $P_+$
selects the $m\ge0$ sector (outer state):
\[
  P_+ = \sum_{n,\,m\ge0} |n,m\rangle\langle n,m|.
\]
Action on the two chiral states:
\[
  P_+|1,{+}1\rangle = |1,{+}1\rangle \quad\text{(outer state survives)},
  \qquad
  P_+|1,{-}1\rangle = 0 \quad\text{(inner state lies in the kernel)}.
\]
After projection the inner state is $|0\rangle$ -- a null vector.
The weak coupling:
\[
  g_L = \langle 1,{+}1|H_\text{weak}|1,{+}1\rangle \neq 0,
  \qquad
  g_R = \langle 0|H_\text{weak}|0\rangle = 0.
\]
\textbf{This is not a sign argument ($g_R=-g_L$) but a kernel argument: the
inner state lies in the kernel of $P_+$ and has no physical existence after
projection.} $g_R=0$ follows algebraically.

\textbf{[S]} \textbf{What remains open: why $P_+$ and not $P_-$?}
The projector $P_+$ corresponds to the outer side ($K>0$, positive orientation,
electric flux inward). $P_-$ would be the inner-state projector. That $P_+$ is
the physical projector follows from orientation arguments (Gaussian curvature,
field rotation $\nabla\times\mathbf{E}_\text{field}$) -- but not from a
derivable symmetry breaking. This is the one remaining open edge.

\section{Symbol Glossary}

Symbols used throughout the entire A-series are explained in A010.
Specific to this document:

\begin{center}
\renewcommand{\arraystretch}{1.3}
{\small\begin{tabular}{lp{9cm}}
\toprule
Symbol & Meaning \\
\midrule
$(n_\theta,n_\phi)=(1,1)$ & Winding numbers for spin-$\tfrac12$ on $T^4$ \\
$K$ & Gaussian curvature; outside $K>0$, inside $K<0$ \\
$\Phi=n\cdot h/e$ & Quantised flux lines, topologically fixed field direction \\
$g_L, g_R$ & Left- and right-handed weak coupling \\
Type-I projection & Information-preserving decompactification (A090) \\
Type-II loss & Non-reversible spatial projection, information loss \\
\bottomrule
\end{tabular}}
\renewcommand{\arraystretch}{1.0}
\end{center}

\section{Verification Script}

\begin{itemize}[leftmargin=2em,itemsep=2pt,topsep=2pt]
  \item Topological consistency: winding $(1,1)$, 720° symmetry, Gaussian
    curvature inside/outside, phase difference $\pm\xi/2$:
    \texttt{python/A\_Serie\_Skripte/a095\_torus\_chiralitaet.py}
\end{itemize}

\section{Open Questions}

\textbf{$g_R=0$ is now at [B] level:} given $P_+$ as the physical projector,
$g_R=0$ follows algebraically from $P_+|1,{-}1\rangle=0$.

\textbf{What remains [S]: why $P_+$ is the physical projector.}
The orientation argument (Gaussian curvature, field rotation) is geometrically
plausible, but not derived from a deeper symmetry breaking. This is the
remaining open edge -- and it is precisely named.

\textbf{Connection to the gauge sector (A190) missing.} The chiral projection
explains the $SU(2)_L$ coupling qualitatively; whether it reproduces the full
electroweak gauge structure is not shown.

\section{Supersedes}

This document systematises the inside/outside chirality mechanism from the
legacy corpus (particle mass torus model, Dirac winding structure, projection
chain framework). It supplements A263 (Dirac, half-integrality), A090
(projection chain, type-I/II), and A264 (asymmetry, CP) with the missing link.

\section{Use of AI Tools}

This document was created with the assistance of AI tools.
Responsibility for content and errors rests with the author.
Full wording of the rules in A010.
